\documentclass[11pt, a4paper]{article}

% ============================================================
% PACKAGES
% ============================================================
\usepackage[utf8]{inputenc}
\usepackage[T1]{fontenc}
\usepackage{lmodern}
\usepackage[margin=1in]{geometry}
\usepackage{graphicx}
\usepackage{hyperref}
\usepackage{listings}
\usepackage{xcolor}
\usepackage{enumitem}
\usepackage{booktabs}
\usepackage{fancyhdr}
\usepackage{titlesec}
\usepackage{parskip}

% ============================================================
% STYLING
% ============================================================

% Link colors
\hypersetup{
    colorlinks=true,
    linkcolor=blue!70!black,
    urlcolor=blue!70!black,
    citecolor=blue!70!black
}

% Code listing style
\lstdefinestyle{codestyle}{
    backgroundcolor=\color{gray!10},
    basicstyle=\ttfamily\small,
    breaklines=true,
    frame=single,
    rulecolor=\color{gray!30},
    numbers=left,
    numberstyle=\tiny\color{gray},
    tabsize=4,
    showstringspaces=false,
    keywordstyle=\color{blue!70!black}\bfseries,
    commentstyle=\color{green!50!black},
    stringstyle=\color{red!60!black}
}
\lstset{style=codestyle}

% Header and footer
\pagestyle{fancy}
\fancyhf{}
\fancyhead[L]{\small\textit{Design Document}}
\fancyhead[R]{\small\textit{Project Name}}
\fancyfoot[C]{\thepage}
\renewcommand{\headrulewidth}{0.4pt}

% Section styling
\titleformat{\section}{\large\bfseries}{}{0em}{}[\vspace{-0.5em}\rule{\textwidth}{0.4pt}]

% ============================================================
% METADATA — FILL THESE IN
% ============================================================
\newcommand{\doctitle}{Feature Name: Design Document}
\newcommand{\docauthor}{Your Name}
\newcommand{\docreviewer}{Reviewer Name(s)}
\newcommand{\docstatus}{Draft}           % Draft | In Review | Approved
\newcommand{\docversion}{0.1}

% ============================================================
% DOCUMENT START
% ============================================================
\begin{document}

% --------------------------------------------------
% Title block
% --------------------------------------------------
\begin{center}
    {\LARGE\bfseries \doctitle \par}
    \vspace{1em}
    \begin{tabular}{@{} l l @{}}
        \textbf{Author:}   & \docauthor \\
        \textbf{Reviewer:} & \docreviewer \\
        \textbf{Status:}   & \docstatus \\
        \textbf{Version:}  & \docversion \\
        \textbf{Date:}     & \today \\
    \end{tabular}
\end{center}

\vspace{1em}
\tableofcontents
\newpage

% --------------------------------------------------
% 1. Overview
% --------------------------------------------------
\section{Overview}

Provide a brief, high-level summary of the feature or system being designed. A reader should understand the ``what'' and ``why'' within a few sentences.

\textit{Example: This document describes the design for a real-time notification service that delivers push notifications to mobile and web clients with at-least-once delivery guarantees.}

% --------------------------------------------------
% 2. Context and Motivation
% --------------------------------------------------
\section{Context and Motivation}

Explain the background. Why is this work needed? What problem does it solve? Include links to relevant docs, tickets, or prior art.

\begin{itemize}
    \item What is the current state of the system?
    \item What pain points or limitations exist today?
    \item What is the business or user impact?
\end{itemize}

% --------------------------------------------------
% 3. Goals and Non-Goals
% --------------------------------------------------
\section{Goals and Non-Goals}

\subsection*{Goals}
\begin{itemize}
    \item Goal 1: Clearly state what this design aims to achieve.
    \item Goal 2: Be specific and measurable where possible.
\end{itemize}

\subsection*{Non-Goals}
\begin{itemize}
    \item Non-Goal 1: Explicitly call out what is out of scope.
    \item Non-Goal 2: This prevents scope creep and sets expectations.
\end{itemize}

% --------------------------------------------------
% 4. Proposed Design
% --------------------------------------------------
\section{Proposed Design}

This is the core of the document. Describe the solution in detail.

\subsection{Architecture Overview}

Describe the high-level architecture. Include a diagram if helpful.

% Uncomment to include a diagram:
% \begin{figure}[h]
%     \centering
%     \includegraphics[width=0.8\textwidth]{architecture-diagram.png}
%     \caption{System architecture overview.}
% \end{figure}

\subsection{Data Model}

Define the key data structures, database schemas, or models.

\begin{table}[h]
    \centering
    \begin{tabular}{@{} l l l @{}}
        \toprule
        \textbf{Field} & \textbf{Type} & \textbf{Description} \\
        \midrule
        id             & UUID          & Primary key \\
        user\_id       & UUID          & Foreign key to users table \\
        status         & ENUM          & pending, active, archived \\
        created\_at    & TIMESTAMP     & Record creation time \\
        \bottomrule
    \end{tabular}
    \caption{Example schema for the notifications table.}
\end{table}

\subsection{API Design}

Define the key endpoints or interfaces.

\begin{lstlisting}[language=Python, caption={Example endpoint}]
@app.post("/api/v1/notifications")
async def create_notification(payload: NotificationRequest):
    """Create and dispatch a new notification."""
    notification = NotificationService.create(payload)
    await NotificationService.dispatch(notification)
    return {"id": notification.id, "status": "pending"}
\end{lstlisting}

\subsection{Key Flows}

Walk through the most important user or system flows step by step.

\begin{enumerate}
    \item User triggers an action in the app.
    \item The API server creates a notification record.
    \item The dispatcher places the message on a queue.
    \item Workers consume the message and deliver it to the client.
\end{enumerate}

% --------------------------------------------------
% 5. Alternatives Considered
% --------------------------------------------------
\section{Alternatives Considered}

Describe other approaches you evaluated and why you chose the proposed design over them. This shows due diligence and helps reviewers understand trade-offs.

\begin{description}
    \item[Option A: Polling-based approach] \hfill \\
        Simpler to implement but adds latency and increases server load at scale.
    \item[Option B: Third-party service (e.g., Firebase)] \hfill \\
        Lower maintenance burden but introduces a vendor dependency and limits customization.
\end{description}

% --------------------------------------------------
% 6. Trade-offs and Risks
% --------------------------------------------------
\section{Trade-offs and Risks}

Be honest about the downsides and uncertainties.

\begin{itemize}
    \item \textbf{Complexity:} The proposed design adds a message queue dependency.
    \item \textbf{Latency:} End-to-end delivery time depends on queue throughput.
    \item \textbf{Risk:} If the queue becomes a bottleneck, notifications may be delayed.
    \item \textbf{Mitigation:} We will implement dead-letter queues and monitoring alerts.
\end{itemize}

% --------------------------------------------------
% 7. Milestones and Timeline
% --------------------------------------------------
\section{Milestones and Timeline}

\begin{table}[h]
    \centering
    \begin{tabular}{@{} l l l @{}}
        \toprule
        \textbf{Milestone} & \textbf{Target Date} & \textbf{Owner} \\
        \midrule
        Design review complete    & YYYY-MM-DD & Author \\
        Core implementation       & YYYY-MM-DD & Team \\
        Integration testing       & YYYY-MM-DD & Team \\
        Staged rollout            & YYYY-MM-DD & Team \\
        \bottomrule
    \end{tabular}
\end{table}

% --------------------------------------------------
% 8. Open Questions
% --------------------------------------------------
\section{Open Questions}

List unresolved decisions or items that need further discussion.

\begin{enumerate}
    \item Should we support batched notifications in v1 or defer to v2?
    \item What is the expected peak throughput, and do we need to load test before launch?
    \item Who owns the on-call rotation for this service?
\end{enumerate}

% --------------------------------------------------
% 9. References
% --------------------------------------------------
\section{References}

\begin{itemize}
    \item \href{https://example.com}{Link to related design doc or RFC}
    \item \href{https://example.com}{Link to relevant ticket or epic}
    \item \href{https://example.com}{Link to external resources or prior art}
\end{itemize}

\end{document}